(1) 叠加场平方给出 $P(t)=P_s+P_{\rm LO}+2\sqrt{P_sP_{\rm LO}}\cos(\Delta\omega t+\phi)$。因此 $I_{\rm dc}=\mathcal R(P_s+P_{\rm LO})$，拍频电流幅值 $I_b=2\mathcal R\sqrt{P_sP_{\rm LO}}$；假设模式完全重合且探测带宽覆盖拍频。
(2) 对Poisson光电子计数，单边电流噪声谱 $S_i=2eI_{\rm dc}$（单位 $\mathrm{A^2/Hz}$），有效噪声带宽 $B$ 中 $\langle i_n^2\rangle=2eI_{\rm dc}B$。本振光自身也贡献散粒噪声，不可只计算信号光的噪声。
(3) 正弦拍频均方电流为 $I_b^2/2=2\mathcal R^2P_sP_{\rm LO}$，故 $\mathrm{SNR}=\mathcal R P_sP_{\rm LO}/[e(P_s+P_{\rm LO})B]$。当 $P_{\rm LO}\gg P_s$，$\mathrm{SNR}\to\mathcal R P_s/(eB)$，不再随本振功率增长。
(4) 若 $\mathcal R=\eta e/(\hbar\omega_0)$，强本振极限为 $\mathrm{SNR}=\eta P_s/(\hbar\omega_0B)$，等于单位有效积分时间内探测的信号光子数尺度。进一步提高本振仅同比增加拍频幅度和本振散粒噪声，还可能引入饱和及技术噪声。
